% Part: many-valued-logic
% Chapter: syntax-and-semantics
% Section: introduction

\documentclass[../../../include/open-logic-section]{subfiles}

\begin{document}

\olfileid{mvl}{syn}{int}

\olsection{پېژندنه}

په کلاسيک منطق کښې له هغو !!{formula}s سره کار
کوو چې له !!{propositional variable}s څخه د
قضيوي رابطونو $\lnot$، $\land$، $\lor$، $\lif$
او $\liff$ په وسيله جوړېږي۔ د کلاسيک منطق د
معنٰی پوهنې د تعريف دپاره د رښتياوالي دوه قيمتونه $\True$
او $\False$ کاروو۔ !!{propositional variable}s
په !!a{valuation}~$\pAssign{v}$ کښې تفسيروو؛
دا هر !!{propositional variable}s ته له دغو دوو څخه يو قيمت،
$\True$ يا $\False$، ورکوي۔ بيا هره
!!{valuation} د هر !!{formula}~$!A$ د رښتياوالي
قيمت $\pValue{v}(!A)$ ټاکي؛ او !!^a{formula}
په !!a{valuation}~$\pAssign{v}$ کښې هله پوره
کېږي، $\pSat{v}{!A}$، چې $\pValue{v}(!A) =
\True$ وي۔

څو ارزښته منطقونه د کلاسيک دوه ارزښته منطق
عموميتونه دي: د $\True$ او $\False$ تر څنګ
د رښتياوالي نور قيمتونه هم مني۔ نو په څو ارزښته منطق کښې
!!a{valuation}~$\pAssign{v}$ داسې تابعه ده چې
هر !!{propositional variable}~$p$ ته د رښتياوالي
د ممکنه قيمتونو له ټولګې څخه يو قيمت ورکوي۔
د رښتياوالي د منل شوو قيمتونو سټ ته به عموماً~$V$
وايو۔ کلاسيک منطق هم يو څو ارزښته منطق دے چې
$V = \{\True, \False\}$ لري؛ د فارمول د رښتياوالي
قيمت~$\pValue{v}(!A)$ د رابطونو د پېژندل شوو
ځانګړنو د رښتياوالي جدولونو له مخې محاسبه کېږي۔

کله چې د رښتياوالي نور قيمتونه ور زيات کړو، د هغو
رابطونو د $\pValue{v}(!A)$ د محاسبې دپاره چې
«او»، «يا»، «نۀ» او «که---نو» ئې لولو، له يوې
څخه زياتې طبيعي لارې شته۔ نو يو څو ارزښته منطق
يوازې د رښتياوالي د قيمتونو سټ نۀ ټاکي؛ د هر رابط
دپاره ټاکل شوې \emph{د صدق تابعې} هم ئې ټاکي۔
خو چې د يوه څو ارزښته منطق~$\Log L$ دپاره دغه
تابعې وټاکل شي، د رښتياوالي قيمت
$\pValue{v}(!A)[\Log L]$ د ارزښت ټاکنې له مخې
په يوازيني ډول معلومېږي، لکه په کلاسيک منطق کښې۔
څو ارزښته منطقونه هم، د کلاسيک منطق په څېر،
\emph{صدق‌تابع} دي۔

له دې معنايي بنسټونو څخه د تاوتولوژۍ،
استلزام او د صدق وړتيا د معنايي مفهومونو سيالې
بڼې تعريفولے شو۔ په کلاسيک منطق کښې
!!a{formula} هله تاوتولوژي وي چې د هرې
$\pAssign{v}$ دپاره ئې د رښتياوالي قيمت
$\pValue{v}(!A) = \True$ وي۔ په څو ارزښته
منطق کښې دا خبره لږ عمومي کوو۔ لومړے، لازمه نۀ
ده چې د رښتياوالي د قيمتونو سټ~$V$ پخپله $\True$
ولري۔ د بېلګې په توګه، ځينې څو ارزښته منطقونه
عددونه کاروي: له $0$ څخه تر~$1$ پورې ټول ناطق
عددونه د رښتياوالي د قيمتونو سټ ګڼي۔ په دې حالت کښې
$1$ عموماً د~$\True$ رول لوبوي؛ په ځينو نورو
منطقونو کښې څو قيمتونه دا رول لري۔ نو له هر څو
ارزښته منطق څخه غواړو چې د \emph{ټاکل شوو
قيمتونو} يو سټ~$V^+$ ولري۔ بيا ويلے شو چې
!!a{formula} په !!a{valuation}~$\pAssign{v}$
کښې هله پوره کېږي، $\pSat{v}{!A}[\Log L]$،
چې $\pValue{v}(!A)[\Log L] \in V^+$ وي۔
!!^a{formula}~$!A$ د دغه منطق تاوتولوژي،
$\Entails[\Log L] !A$، هله ده چې د هرې
$\pAssign{v}$ دپاره $\pValue{v}(!A) \in V^+$
وي۔ په پاې کښې وايو چې د !!{formula}s سټ
د $!A$ استلزام کوي، $\Gamma \Entails[\Log L] !A$،
که هره هغه !!{valuation} چې په~$\Gamma$ کښې
ټول !!{formula}s پوره کوي، $!A$ هم پوره کوي۔

\end{document}
